\documentclass[11pt]{article}

\usepackage[margin=1in]{geometry}
\usepackage{amsmath,amssymb,amsfonts,bm}
\usepackage{booktabs,longtable,array}
\usepackage{graphicx}
\usepackage{xcolor}
\usepackage{hyperref}
\usepackage{enumitem}
\usepackage{algorithm}
\usepackage{algorithmic}

\hypersetup{
  colorlinks=true,
  linkcolor=blue!55!black,
  citecolor=blue!55!black,
  urlcolor=blue!55!black
}

\newcommand{\R}{\mathbb{R}}
\newcommand{\C}{\mathbb{C}}
\newcommand{\dd}{\mathrm{d}}
\newcommand{\diag}{\operatorname{diag}}
\newcommand{\rank}{\operatorname{rank}}
\newcommand{\spec}{\operatorname{spec}}
\newcommand{\Real}{\operatorname{Re}}
\newcommand{\Imag}{\operatorname{Im}}
\newcommand{\norm}[1]{\left\lVert #1 \right\rVert}
\newcommand{\Lt}{\widetilde L}
\newcommand{\Dt}{\widetilde D}
\newcommand{\Sigmae}{\Sigma}
\newcommand{\zetaReq}{\zeta_{\rm req}}

\newtheorem{theorem}{Theorem}
\newtheorem{proposition}{Proposition}
\newtheorem{conjecture}{Conjecture}
\newtheorem{definition}{Definition}
\newtheorem{remark}{Remark}

\title{\bfseries Internal Unified Theory Note:\\
Damping-Operator Hierarchy, Rational Self-Energy, and Flow-Feasible Retrofitting in IBR Grids}
\author{Project working document}
\date{June 2026}

\begin{document}
\maketitle

\begin{abstract}
This document is the internal map of the project. It is not a submission draft.
It organizes the verified results, failed hypotheses, and open gates under one
operator hierarchy. The central thesis is that the transition from synchronous
grids to inverter-rich grids is a descent through a damping-operator hierarchy:
proportional viscous damping, non-proportional viscous damping, and rational
frequency-dependent damping induced by controller self-energy. Each descent
breaks a classical guarantee and creates a planning risk. Level 0 gives normal
modes and topology-invariant decay under proportional damping. Level 1 breaks
classical modal separation; heterogeneous $D/M$ creates modal coupling, dynamic
Braess-like reversals, and non-monotone inertia effects. Level 2 condenses
controller states into a rational self-energy
\[
  \Sigma(s)=D_0+C(sI-A_{cc})^{-1}B,
\]
which yields a nonlinear eigenvalue problem. This bridge was validated against
full ANDES eigensolves on IEEE 39-bus cases, including a 13.7417 Hz PLL-family
critical mode in Mix60. The document also integrates the newer resolvent-margin
and coordinated-retrofitting gates. Those gates support a secondary but useful
claim: modal damping-ratio and protected-resolvent certificates rank actions
differently across several protection weightings, and equal-cost packages
$w+D$ and $w+M$ often beat the best single lever under the protected resolvent.
The mechanism of that certificate discrepancy is still open; simple global
non-normality and commutator descriptors correlated weakly. Every claim below
is labeled by validation status.
\end{abstract}

\tableofcontents
\newpage

\section{Executive Thesis}

The project can be read as one question:
\[
  \text{What object preserves damping-margin mechanisms after IBR controls are condensed?}
\]
The answer changed as the work progressed. A scalar $(L,M,D)$ bridge is useful
for intuition but not enough for detailed IBR dynamics. The correct reduced
object for the ANDES cases is a rational Schur self-energy or, equivalently,
a retained-state nonlinear eigenvalue problem (NEP). Once that bridge exists,
we can ask smaller planning questions: which line is weak, which node is weak,
which damping/inertia/topology action moves the margin, and whether packages of
actions beat single levers under flow feasibility.

\subsection{The operator hierarchy}

The useful organizing device is a three-level hierarchy of damping operators.
Each level has a mathematical owner in the literature, a guarantee, and a
planning risk when the grid leaves that level.

\begin{center}
\renewcommand{\arraystretch}{1.25}
\begin{tabular}{p{0.14\textwidth}p{0.26\textwidth}p{0.27\textwidth}p{0.19\textwidth}}
\toprule
Level & Operator & Classical guarantee & Planning risk when violated\\
\midrule
0 & Proportional viscous damping, $D=\gamma M$ or Rayleigh/Caughey-class damping
& Normal-mode separation; common decay for underdamped proportional modes; topology moves frequency more than decay
& Frequency-only planning can look valid, but only inside the proportional floor\\
1 & Constant non-proportional viscous damping, $[\Lt,\Dt]\neq0$
& Modal coordinates are no longer decoupled; off-diagonal modal damping moves poles
& Line reinforcement or inertia can reduce damping ratio; Fiedler/frequency rankings can fail\\
2 & Rational/nonviscous damping, $\Sigma(s)=D_0+C(sI-A_{cc})^{-1}B$
& Constant-coefficient QEP is no longer the reduced object; controller poles can become critical
& Static bridge and fixed-point screens can miss control-family critical modes\\
\bottomrule
\end{tabular}
\end{center}

This hierarchy is not a claim that Rayleigh damping, non-proportional modal
perturbation, nonviscous damping, Schur complements, or contour eigensolvers are
new. They are not. The project contribution, if it survives adversarial review,
is the power-system formulation and validation: treating IBR grids as active
nonviscous networked structures and using a controller-aware rational bridge to
classify network-family and control-family margin mechanisms.

\subsection{The material-active candidate framing}

The most compact physical sentence is:

\begin{quote}
An IBR-dominated grid behaves like a viscoelastic structure made of active
material: its controller self-energy can have oscillatory, band-active poles,
so nonviscous modes can become the critical small-signal modes.
\end{quote}

This is a candidate framing, not a closed novelty claim. Passive viscoelastic
models have a long literature, including hereditary kernels and nonviscous
damping models. The power-system observation is that a PLL or inverter-control
kernel is not a passive relaxation kernel in the same sense. In Mix60, the
critical pole is a 13.7417 Hz control-family mode with $\pi_c=0.853$, close to
a PLL-dominated condensed-control pole. That is the empirical anchor. The
adversarial literature search for ``active viscoelastic structures,'' active
nonviscous kernels, and nonviscous damping in power grids remains open.

\subsection{What papers can come out}

\begin{itemize}[leftmargin=*]
\item \textbf{Paper 1: the controller-aware bridge.} This is the strongest
validated paper. It reports the scalar bridge failure, the Schur NEP bridge,
Beyn contour extraction, network/control-family classification, Mix60's
control-family critical mode, and the resonant Braess mechanism with negative
control.
\item \textbf{Paper 2: the hierarchy and passivity gate.} This is conditional.
It becomes strong if a passivity/active-kernel test separates cases where the
static bridge is safe from cases where Level 2 is required.
\item \textbf{Paper 3: flow-feasible resolvent retrofitting.} This is a
conference-scale candidate. The Gate 1 and Gamma gates show a real phenomenon:
protected-resolvent and modal certificates recommend different actions, and
$w+D$ or $w+M$ packages often beat single levers under the protected resolvent.
The missing piece is a convincing local mechanism descriptor.
\item \textbf{Thesis material.} The weak-node/weak-link maps, substitutability
spectrum, inertia reversals, and graph-configuration optimization are useful
thesis chapters. Several are formulated or surrogate-verified rather than full
ANDES-validated.
\end{itemize}

\section{Validation Ledger}

Table~\ref{tab:ledger} is the most important table in the document. It separates
validated bridge results from formulated planning ideas and dead ends.

\begin{longtable}{p{0.20\textwidth}p{0.25\textwidth}p{0.24\textwidth}p{0.20\textwidth}}
\caption{Claim ledger. Status terms are intentionally blunt.}\label{tab:ledger}\\
\toprule
Claim & Evidence & Control or comparison & Status\\
\midrule
\endfirsthead
\toprule
Claim & Evidence & Control or comparison & Status\\
\midrule
\endhead
Scalar $(L,M,D)$ bridge is insufficient for detailed ANDES damping
& In IEEE39 no-PSS/base/Mix60, local scalar damping from GENROU.D is zero while modal damping exists; Phase-0D reports scalar $D$ norm zero and nonzero Schur damping diagnostics.
& Full ANDES eigensolve and Schur diagnostics.
& ANDES-validated failure mode\\
\midrule
Controller-aware Schur NEP bridge recovers critical poles
& Phase-0D: no-PSS, base, Mix60 pass. Worst reported errors: Mix60 frequency $8.475\times10^{-11}\%$, damping $2.479\times10^{-9}\%$.
& Full ANDES eigensolve; contour argument counts; non-circular polish.
& ANDES-validated\\
\midrule
Mix60 critical mode is control-family
& Pole $-10.7514+j86.3419$, $f=13.7417$ Hz, $\zeta=0.12357$, $\pi_c=0.853$; nearest $A_{cc}$ pole distance 0.215; top reconstructed states PLL1/REGF1.
& Scalar bridge and fixed-point Schur reduction missed or misidentified the mode.
& ANDES-validated mechanism anchor\\
\midrule
Resonant Braess mechanism exists in benchmark
& Phase Braess-NEP: 40/348 strict resonant records. Clearest case Line\_42 (20--34), self-energy term $C$ dominates with control denominator fraction 0.9594 and median analytic-vs-NEP smooth error $8.08\times10^{-4}$.
& Requires network-family damping decrease, closer $A_{cc}$ pole, PLL-dominated nearest pole, and dominant self-energy term; soft-GFL negative control removes resonances.
& ANDES/NEP mechanism validated, minority and local\\
\midrule
Second-order modal correction improves diagonal screen in weak/mild regimes
& Synthetic and reduced-PLL audits: weak surrogate improvement from diagonal to second order; reduced PLL reported $2.16\%\to0.09\%$ median in prior cross-check.
& Diagonal screen and reduced QEP fallback.
& Numerically verified outside full ANDES; estimator validation pending\\
\midrule
Second-order correction always beats large reduced QEP
& Refuted by cross-checks; QEP $r=20$ can beat or match second order in tails.
& Reduced QEP size sweep.
& KILLED\\
\midrule
Inertia can reduce damping margin
& Fixed-$D_i/M_i$ surrogate audit: 5447/16229 perturbations reduce both diagonal and full-QEP margin. ANDES GENROU.M gate: no-PSS/base have 5/10 negative GENROU.M perturbations.
& Fixed-$D_i$ versus fixed-$D_i/M_i$ control; full ANDES eigensolve for GENROU.M cases.
& Valid for mechanical inertia checks; virtual inertia GFM blocked\\
\midrule
Virtual-inertia assignment works over controller-aware IBR bridge
& REGF1 sheets expose no clear $M$, $H$, or $T_j$ knob. Mix60 response to retained GENROU.M is weak signal.
& Phase inertia bridge gate.
& BLOCKED; do not claim\\
\midrule
Certificate discrepancy $S_\zeta$ vs $S_g$
& Gate 1: ranking agreement 28.3\%. Gamma gate: physical profiles have agreement 21--29\%.
& Four physical Gamma choices and eight random diagonal profiles.
& Robust phenomenon, mechanism weak\\
\midrule
Coordinated packages beat single levers at equal cost
& Gate 1/Gamma: under $S_g$, $w+D$ has $CV<0$ in 72.1--73.8\%, $w+M$ in 67.2--73.8\%; $D+M$ not robust.
& Equal-cost budget, DC flow feasibility, same seeds, Gamma profiles.
& Reduced-model gate supports $w+D$ and $w+M$ only\\
\midrule
Nonlinear synergy of levers
& SynNorm $<-0.10$ never exceeds 50\% in Gate 1.
& Three eps sizes and all lever pairs.
& KILLED as headline\\
\midrule
Cross inertia-control term is empirically frequent in ANDES
& Phase E0 found 0/13 measurable singleton cases.
& ANDES bridge checks.
& KILLED as empirical headline; remains local theory\\
\bottomrule
\end{longtable}

\section{Level 0: The Classical Floor}

The starting point is the second-order model
\begin{equation}
  M\ddot x+D\dot x+Lx=0,
  \label{eq:level0-second-order}
\end{equation}
where $M\succ0$ and $L=L^T\succeq0$. Inertia-normalized coordinates
$y=M^{1/2}x$ give
\begin{equation}
  \ddot y+\Dt\dot y+\Lt y=0,\qquad
  \Lt=M^{-1/2}LM^{-1/2},\quad
  \Dt=M^{-1/2}DM^{-1/2}.
  \label{eq:normalized-coords}
\end{equation}
If $D=\gamma M$, then $\Dt=\gamma I$ and every eigenvector of $\Lt$ is also
an eigenvector of $\Dt$. Writing $\Lt q_k=\mu_kq_k$, each nonzero mode obeys
\begin{equation}
  \ddot z_k+\gamma\dot z_k+\mu_k z_k=0,
  \label{eq:level0-scalar}
\end{equation}
with poles
\begin{equation}
  s_{k,\pm}=\frac{-\gamma\pm\sqrt{\gamma^2-4\mu_k}}{2}.
  \label{eq:level0-poles}
\end{equation}
For underdamped modes, $\mu_k>\gamma^2/4$, the real part is exactly
\begin{equation}
  \Real(s_{k,\pm})=-\gamma/2.
  \label{eq:rigid-decay}
\end{equation}

\begin{proposition}[Topology-invariant decay at the proportional floor]
Under $D=\gamma M$, every mode that remains underdamped has decay rate
$-\gamma/2$, independent of topology.
\end{proposition}

\noindent This is a floor theorem, not a project contribution. It is the same
normal-mode logic behind classical proportional damping and Rayleigh damping.
It explains why many graph-frequency intuitions feel safe in homogeneous
systems. The escape channel is also clear: topology can move $\mu_k$ across
the overdamped/underdamped boundary. The rigorous statement is not that
topology never matters. It is that decay is invariant for modes that stay in
the underdamped proportional class.

The project validation notes recorded numerical deviations at the level of
$5.8\times10^{-16}$ to $1.4\times10^{-16}$ in proportional tests. Those numbers
are useful sanity checks, not new theory.

\section{Level 1: Non-Proportional Viscous Damping}

Level 1 begins when $\Dt$ and $\Lt$ no longer share eigenvectors. In modal
coordinates of $\Lt$,
\begin{equation}
  \ddot z+\Gamma\dot z+\Lambda z=0,\qquad
  \Gamma=Q^T\Dt Q,\quad \Lambda=\diag(\nu_k).
  \label{eq:level1-modal}
\end{equation}
Write
\begin{equation}
  \Gamma=\Delta+E,\qquad \Delta=\diag(\delta_k),\qquad E_{kk}=0.
  \label{eq:level1-offdiag}
\end{equation}
The diagonal screen uses scalar modal polynomials
\begin{equation}
  p_k(s)=s^2+\delta_ks+\nu_k.
  \label{eq:scalar-modal-polynomial}
\end{equation}
The actual QEP is coupled by $sE$. This is the first place where topology,
damping, and inertia can interact in non-intuitive ways.

\subsection{QEP sensitivities}

For the QEP
\begin{equation}
  P(s)v=(s^2M+sD+L)v=0,
  \label{eq:qep-general}
\end{equation}
assume a simple eigenvalue and use the complex-symmetric bilinear product
$v^T(\cdot)v$ when $M,D,L$ are real symmetric. Define
\begin{equation}
  \xi=v^T(2sM+D)v.
  \label{eq:qep-xi}
\end{equation}
Then the local parameter sensitivities are
\begin{align}
  \frac{\partial s}{\partial M_i}
  &= -\frac{s^2 v_i^2}{\xi}, \label{eq:qep-sens-M}\\
  \frac{\partial s}{\partial D_i}
  &= -\frac{s v_i^2}{\xi}, \label{eq:qep-sens-D}\\
  \frac{\partial s}{\partial w_e}
  &= -\frac{(b_e^Tv)^2}{\xi}, \label{eq:qep-sens-w}
\end{align}
where $b_e$ is the signed incidence vector of line $e$ in the selected
coordinates. These formulas are standard perturbation theory. In the project
they were used as a common local currency for topology, damping, and inertia.
They were checked against central differences with relative errors around
$10^{-8}$ in the three channels in the associated validation notes.

\subsection{Why inertia can be deceptive}

At the proportional point $D=\gamma M$, holding $D$ fixed while perturbing $M_i$
breaks proportionality. That direction is not the same as increasing both
$M_i$ and $D_i$ proportionally. For an underdamped proportional mode, the
directional derivative of the real part can move the pole right; in the
project notes the exact local rate was written as a positive proportional
quantity of the modal participation. The correct interpretation is damping
dilution: if damping is held fixed and inertia is added, the ratio between
energy storage and dissipation changes.

The stronger phrase ``inertia is universally Braessian'' is false and must not
be used. Monte Carlo and fixed-ratio controls rejected it. The defensible
statement is narrower:

\begin{quote}
Under heterogeneous damping, increasing inertia frequently but not universally
reduces damping margin; at the proportional point, the harmful direction can be
computed exactly when damping is held fixed.
\end{quote}

The current full-ANDES evidence is also limited. In no-PSS and base IEEE39
checks, retained mechanical GENROU.M perturbations had 5/10 negative margin
effects and the local NEP sensitivity matched full ANDES. In Mix60, the
critical mode was control-family and the retained GENROU.M signal was too weak.
No GFM virtual-inertia claim is available until a physical REGF1 inertia knob
is identified.

\subsection{The line Braess dichotomy}

At Level 0, small-signal Braess-like damping degradation is blocked for
underdamped proportional modes because line actions move modal frequency but
not decay. At Level 1, off-diagonal modal damping allows a line action to move
both the frequency and the decay rate. The damping-ratio derivative is
\begin{equation}
  \frac{\partial \zeta}{\partial \rho}
  = -\frac{1}{|s|}\Real\!\left(\frac{\partial s}{\partial \rho}\right)
    +\frac{\Real(s)}{|s|^3}
      \Real\!\left(\overline{s}\frac{\partial s}{\partial \rho}\right),
  \label{eq:zeta-derivative}
\end{equation}
so a line can improve a graph stiffness proxy while reducing damping ratio.
The project found such reversals in surrogate and NEP audits. The resonant
Level-2 version is stronger and appears below.

\subsection{Second-order modal correction}

Let $s_0$ be a diagonal-screen pole of the critical mode $c$:
\begin{equation}
  p_c(s_0)=s_0^2+\delta_cs_0+\nu_c=0.
  \label{eq:critical-diagonal-pole}
\end{equation}
To second order in $E$, eliminating the noncritical modal coordinates gives
\begin{equation}
  \Delta s_c
  =-\frac{s_0^2}{2s_0+\delta_c}
    \sum_{\ell\ne c}
    \frac{E_{c\ell}E_{\ell c}}
    {s_0^2+\delta_\ell s_0+\nu_\ell}
    +O(\norm{E}^3),
  \label{eq:second-order-correction}
\end{equation}
up to sign convention depending on whether the eliminated term is moved to the
left or right side. The useful object is the denominator:
\begin{equation}
  H_{\ell\ell}(s_0)=\frac{1}{s_0^2+\delta_\ell s_0+\nu_\ell}.
  \label{eq:rational-filter-H}
\end{equation}
Then the correction is a rational graph/modal filter:
\begin{equation}
  \Delta s_c \propto [E\,H(\Lambda,\Delta;s_0)\,E]_{cc}.
  \label{eq:rational-filter-form}
\end{equation}
This is not a new perturbation theorem. Its value in the project is
interpretive: each term says which noncritical mode couples into the critical
margin and how strongly its modal denominator amplifies the effect.

The empirical status is mixed in a useful way. The correction improves the
diagonal screen in weak-to-mild non-proportional regimes and in a reduced PLL
model. It does not replace a large reduced QEP in clustered or strongly coupled
tails. Therefore the estimator is a workflow:
\[
  \text{diagonal screen}\quad\to\quad
  \text{second-order correction}\quad\to\quad
  \text{reduced/full QEP or NEP fallback}.
\]

\section{Level 2: Rational Nonviscous Bridge}

\subsection{Schur self-energy}

Let the full linearized ANDES state matrix be partitioned as
\begin{equation}
  A=
  \begin{bmatrix}
  A_{ee} & A_{ec}\\
  A_{ce} & A_{cc}
  \end{bmatrix},
  \label{eq:schur-partition}
\end{equation}
where $x_e$ are retained network/interface states and $x_c$ are condensed
controller, exciter, stabilizer, PLL, governor, and auxiliary states. The
Schur NEP is
\begin{equation}
  S_e(s)=sI-A_{ee}-A_{ec}(sI-A_{cc})^{-1}A_{ce}.
  \label{eq:schur-nep}
\end{equation}
When the retained coordinates are reduced to a second-order form, the same
structure can be written as
\begin{equation}
  T(s)=s^2M+s\Sigma(s)+\Lt,\qquad
  \Sigma(s)=D_0+C(sI-A_{cc})^{-1}B.
  \label{eq:rational-self-energy}
\end{equation}
This is the Level-2 object. It is a rational nonviscous damping operator.
It is also the correct place to connect power-system IBR controls to the
structural-dynamics literature on hereditary and nonviscous damping.

\subsection{Beyn contour validation}

The Phase-0D gate solved \eqref{eq:schur-nep} by Beyn contour integration.
No fixed-point iteration was used for acceptance. The contour solver first
extracted zeros from contour moments and only afterward compared them with
ANDES poles.

\begin{center}
\small
\begin{tabular}{lcccccc}
\toprule
Case & retained & condensed & pole family & $f$ Hz & $\zeta$ & critical error\\
\midrule
no-PSS & 20 & 110 & network & 1.3705 & 0.15447 & $4.861\times10^{-14}\%$ freq\\
base & 20 & 140 & network & 1.3706 & 0.15443 & $1.62\times10^{-13}\%$ freq\\
Mix60 & 11 & 167 & control & 13.7417 & 0.12357 & $8.475\times10^{-11}\%$ freq\\
\bottomrule
\end{tabular}
\end{center}

For Mix60, the matched pole was
\begin{equation}
  s_{\rm crit}=-10.7514+j86.3419,
  \label{eq:mix60-critical-pole}
\end{equation}
with $\pi_c=0.853$. The nearest condensed-control pole was at distance 0.215
and the reconstructed condensed states were dominated by PLL1 and REGF1
components. The scalar bridge and fixed-point reduction did not capture this
mode. This is the strongest full-ANDES result in the project.

\subsection{Active nonviscous material}

Passive nonviscous damping kernels in structural dynamics are usually designed
to dissipate energy. Their extra nonviscous modes are not expected to become
the least-damped oscillatory modes in the same way a PLL-family mode can in an
IBR grid. In contrast, the condensed IBR/control self-energy can have poles and
frequency bands that act like active material. A concise way to write the
candidate test is to define a passivity defect:
\begin{equation}
  \Pi_{\rm pass}
  =
  \max_{\omega\in\Omega}
  \left[-\lambda_{\min}\left(
  \frac{\Sigma(j\omega)+\Sigma(j\omega)^*}{2}
  \right)\right]_+ .
  \label{eq:passivity-defect}
\end{equation}
The conjectured rule is:

\begin{conjecture}[Passive-kernel screening safety]
If $\Pi_{\rm pass}$ is small on the relevant band and the Level-1 Caughey defect
\begin{equation}
  \chi_{\rm CK}=\frac{\norm{\Lt\Dt-\Dt\Lt}_F}{\norm{\Lt}_F\norm{\Dt}_F}
  \label{eq:caughey-defect}
\end{equation}
is small, then Level-0/1 screens are reliable for network-family modes. If
$\Pi_{\rm pass}$ is large, control-family modes can become critical and the
full rational bridge is required.
\end{conjecture}

This is not yet a theorem. The Gamma gate already showed that simple global
non-normality and global $\chi$ do not explain certificate discrepancy strongly.
The passivity defect must therefore be tested locally at the critical mode and
controller-pole band, not sold as an established result.

\section{Resonant Braess and Hidden Control Margin}

For a line reinforcement parameter $\rho$, the Schur NEP sensitivity is
\begin{equation}
  \frac{\partial s}{\partial \rho}
  =
  -\frac{y^H S_\rho(s)x}{y^HS_s(s)x},
  \label{eq:schur-sensitivity}
\end{equation}
where $x$ and $y$ are right and left null vectors of $S_e(s)$. The derivative
$S_s(s)$ includes the resolvent-square term
\begin{equation}
  A_{ec}(sI-A_{cc})^{-2}A_{ce}.
  \label{eq:control-resolvent-square}
\end{equation}
This term is the signature of frequency-dependent control self-energy.

The Phase Braess-NEP gate decomposed the damping-ratio movement into retained
network terms and a condensed-control term. A reinforcement was called
resonant only if it satisfied all of the following:
\begin{enumerate}[leftmargin=*]
\item the tracked network-family damping ratio decreased;
\item the tracked mode moved closer to a matched $A_{cc}$ pole;
\item the nearest condensed pole was PLL-dominated;
\item the self-energy term dominated the retained terms;
\item the effect survived finite-difference NEP/ANDES comparison.
\end{enumerate}
The gate found 40 strict resonant records among 348 reversals. The clearest
case was Mix60/no-PSS Line\_42 (20--34), mode rank 2:
\[
\zeta_{\rm margin}:0.123566\to0.123558,\quad
\zeta_{\rm network}:0.133884\to0.133816,
\]
with distance to the matched $A_{cc}$ pole $3.33227\to3.3279$. The self-energy
component was $-6.874\times10^{-3}$, while the retained terms were orders of
magnitude smaller, and the finite-difference derivative was
$-6.860\times10^{-3}$.

This result should be written as a mechanism result. It is not a claim that
every line reinforcement or every IBR grid exhibits resonant Braess behavior.
The soft-GFL negative control moved the system away from the resonance band and
removed the strict resonant cases.

\section{Protected Resolvent Certificates and Coordination}

\subsection{Two margins}

The modal margin used throughout the earlier project is
\begin{equation}
  S_\zeta=\log\frac{\zetaReq}{\zeta_{\rm crit}},
  \label{eq:modal-margin}
\end{equation}
where lower is better and $S_\zeta\le0$ is safe. The protected-resolvent margin is
\begin{equation}
  S_g=\log g_{\rm prot},\qquad
  g_{\rm prot}=\sup_{\omega\in\Omega}
  \sigma_1\!\left(
  \Gamma_y C(j\omega I-A)^{-1}B\Gamma_u
  \right).
  \label{eq:protected-resolvent}
\end{equation}
The Gate 1 baseline used identity protection weights. The Gamma gate repeated
the test with physically motivated and random diagonal protection weights.

\subsection{Bi-orthogonal resolvent sensitivity}

Let $\omega^\star$ be a simple maximizer of \eqref{eq:protected-resolvent}, and
let $u^\star,v^\star$ be the leading left and right singular vectors of
\[
  H^\star=\Gamma_y C(j\omega^\star I-A)^{-1}B\Gamma_u.
\]
Set
\begin{align}
  \phi^{\rm in}&=(j\omega^\star I-A)^{-1}B\Gamma_u v^\star,\\
  \phi^{\rm out}&=(j\omega^\star I-A)^{-H}C^H\Gamma_y^H u^\star .
\end{align}
For a parameter $p$ that changes $A$ but not the protection weights, the local
sensitivity is
\begin{equation}
  \frac{\partial S_g}{\partial p}
  =
  \frac{1}{g_{\rm prot}}
  \Real\left[
  (\phi^{\rm out})^H
  \frac{\partial A}{\partial p}
  \phi^{\rm in}
  \right],
  \label{eq:resolvent-sensitivity}
\end{equation}
provided the peak and singular value are simple. This is standard singular-value
perturbation applied to the resolvent. Its value is that all levers can be
priced in the same margin units.

For the swing first-order matrix
\[
A=\begin{bmatrix}
0&I\\
-M^{-1}L&-M^{-1}D
\end{bmatrix},
\]
write $\phi=(\phi_\theta,\phi_\omega)$. Then local damping at bus $i$ has
\begin{equation}
  \frac{\partial A}{\partial D_i}
  =
  \begin{bmatrix}
  0&0\\
  0&-M_i^{-1}e_ie_i^T
  \end{bmatrix},
  \label{eq:dA-dD}
\end{equation}
so
\begin{equation}
  \frac{\partial S_g}{\partial D_i}
  =
  -\frac{1}{g_{\rm prot}M_i}
  \Real\left[
  \overline{\phi^{\rm out}_{\omega,i}}\,
  \phi^{\rm in}_{\omega,i}
  \right].
  \label{eq:damping-resolvent-placement}
\end{equation}
This is the corrected version of the Gemini suggestion. It is not generally
$-|\phi_i^{\rm out}|^2$ unless the system is normal or the input and output
spatial vectors align. The important point is the bi-orthogonal product. In a
non-normal grid, the forcing vector and response vector can be spatially
different. That is the mathematical source of spatial decoupling.

For a line $e=(i,j)$ with incidence $b_e$, $L$ changes by $b_eb_e^T$, giving
\begin{equation}
  \frac{\partial A}{\partial w_e}
  =
  \begin{bmatrix}
  0&0\\
  -M^{-1}b_eb_e^T&0
  \end{bmatrix},
  \label{eq:dA-dw}
\end{equation}
and therefore
\begin{equation}
  \frac{\partial S_g}{\partial w_e}
  =
  -\frac{1}{g_{\rm prot}}
  \Real\left[
  (\phi^{\rm out}_\omega)^H
  M^{-1}b_e b_e^T
  \phi^{\rm in}_\theta
  \right].
  \label{eq:line-resolvent-placement}
\end{equation}
This formula explains why a line action is not judged only by flow loading or
graph centrality. It is judged by the overlap between the output response path
and the input forcing path through that edge.

\subsection{Gate results}

Gate 1 tested 60 connected Watts-Strogatz cases with $n=5,\ldots,9$ and two
IEEE topology templates, with DC flow feasibility enforced after topology
changes. Costs were fixed at
\[
  c_w=1,\qquad c_D=2,\qquad c_M=8.
\]
Two subtests were run. Nonlinear synergy was measured by
\begin{equation}
  \operatorname{Syn}_\epsilon(a,b)
  =
  S(p+\epsilon a+\epsilon b)-S(p+\epsilon a)-S(p+\epsilon b)+S(p),
  \label{eq:synergy}
\end{equation}
normalized by $|\Delta S(a)|+|\Delta S(b)|+\eta$. Equal-cost coordination was
measured by
\begin{equation}
  CV(a,b)=\Delta S_{\rm equalcost}(a,b)
  -\min\{\Delta S_{\rm equalcost}(a),\Delta S_{\rm equalcost}(b)\}.
  \label{eq:coordination-value}
\end{equation}
Negative $CV$ means the package beats the best individual lever at the same
total cost.

The nonlinear-synergy headline died. $\operatorname{SynNorm}<-0.10$ never
exceeded 50\% of cases. But equal-cost coordination under $S_g$ survived:
\[
\begin{array}{c|ccc}
\text{Gamma profile} & w+D & w+M & D+M\\
\hline
\text{identity} & 72.1\% & 67.2\% & 51.6\%\\
\text{inertia output} & 73.8\% & 73.8\% & 48.4\%\\
\text{type input low inertia} & 73.8\% & 70.5\% & 48.4\%
\end{array}
\]
Thus $w+D$ and $w+M$ are supported in this reduced-model gate; $D+M$ is not.

The certificate discrepancy also survived Gamma changes. Ranking agreement
between $S_\zeta$ and $S_g$ was 29.0\% for identity, 21.0\% for inertia-output,
22.6\% for type-input-low-inertia, and 29.0\% for the frequency-output control.
Eight random positive diagonal Gamma samples did not remove the discrepancy.

The mechanism is not yet explained. The maximum absolute Spearman correlation
between the tested global descriptors and ranking disagreement was only 0.216.
The correct status is therefore:

\begin{quote}
Certificate discrepancy is robust to tested Gamma choices, but the global
non-normality descriptors and alignment descriptors used so far do not explain
it well.
\end{quote}

\subsection{Flow-feasible resolvent retrofitting algorithm}

The Gemini proposal can be included as an algorithmic candidate, but only after
correcting the sensitivity formulas as above and adding the status labels.

\begin{algorithm}[h]
\caption{Flow-feasible protected-resolvent retrofitting}
\begin{algorithmic}[1]
\REQUIRE Dynamic model $A$, input/output matrices $B,C$, protection weights
$\Gamma_u,\Gamma_y$, DC dispatch $P$, flow limits $\bar f$, action costs.
\STATE Compute $g_{\rm prot}$, $\omega^\star$, and singular vectors
$u^\star,v^\star$ of $\Gamma_y C(j\omega I-A)^{-1}B\Gamma_u$.
\STATE Compute $\phi^{\rm in}$ and $\phi^{\rm out}$.
\STATE For each bus, compute damping and inertia exchange rates using
\eqref{eq:resolvent-sensitivity}. For inertia, use the actual parameter map;
do not substitute a virtual-inertia formula unless the controller exposes one.
\STATE For each line, compute the topology exchange rate
\eqref{eq:line-resolvent-placement} and the DC flow sensitivity/feasibility
penalty.
\STATE Solve a local cost-constrained problem
\[
  \min_z c^Tz\quad
  \text{s.t.}\quad
  S_g+\nabla S_g^Tz\le0,\quad
  |f(w+\Delta w;P)|\le\bar f,\quad z\in\mathcal Z.
\]
\STATE Validate the proposed package by recomputing the full eigensolve and
the protected resolvent. Reject if the local prediction is not matched.
\end{algorithmic}
\end{algorithm}

This algorithm is not yet validated on full ANDES IBR. It is the natural next
step for the protected-resolvent paper.

\section{Weak Nodes, Weak Links, and Substitutability}

The weak-element language should be mode- and certificate-specific. A bus is
not weak in the abstract. It is weak for a margin, a mode family, and a lever
set.

\begin{definition}[Local weak node]
For a margin $S$ and admissible local action set $\mathcal A_i$ at bus $i$,
define
\begin{equation}
  W_i(S)=\max_{a\in\mathcal A_i}
  \left[\nabla S^T a\right]_+ .
  \label{eq:weak-node}
\end{equation}
The node is weak if its best admissible perturbation can increase the risk
margin to first order.
\end{definition}

\begin{definition}[Repairability]
For the same node, define
\begin{equation}
  R_i(S)=\max_{a\in\mathcal A_i,\;c^Ta\le1}
  \left[-\nabla S^T a\right]_+ .
  \label{eq:repairability}
\end{equation}
Weakness and repairability are different. A node can be exposed but expensive
to repair, or easy to repair but not currently exposed.
\end{definition}

For a line $e$, replace $\mathcal A_i$ by admissible line actions and include
the DC feasibility constraint. A line reinforcement that improves a graph
stiffness score but increases $S_\zeta$ or $S_g$ is a dynamic Braess candidate.
It is flow-feasible only if
\begin{equation}
  |f(w+\Delta w;P)|\le \bar f.
  \label{eq:flow-feasible}
\end{equation}

Substitutability compares lever subspaces. In the simplest matrix-geometric
picture, existing-line changes live in
\[
  \mathcal L_{\rm existing}=\operatorname{span}\{b_eb_e^T:e\in E\},
\]
new-line changes live in missing-edge directions, and shunt/inertia-like
changes live in diagonal directions. A mode-specific substitutability score
asks whether the desired pole motion can be achieved by cheap line directions
or requires local equipment. This is a useful planning language. The current
status is formulation plus surrogate audits, not full ANDES validation.

\section{Open Gates}

\begin{enumerate}[leftmargin=*]
\item \textbf{Adversarial literature search for active nonviscous power grids.}
Search terms must include active viscoelastic structures, non-dissipative
nonviscous damping, active material damping kernels, and power-grid NEPs.
\item \textbf{Passivity-defect gate.} Compute \eqref{eq:passivity-defect} for
no-PSS, base, Mix60, soft-GFL controls, and future calibrated IBR cases.
Check whether it predicts when the scalar bridge fails.
\item \textbf{Local discrepancy mechanism.} The Gamma gate rejected global
$d_F(A)$ and global $\chi$ as strong explanations. Test local descriptors:
peak resolvent frequency, singular-vector spatial separation, critical-mode
condition number, and local self-energy denominator fraction.
\item \textbf{Estimator validation over the NEP bridge.} Run diagonal,
second-order, reduced-QEP, and full NEP/ANDES comparisons on calibrated IBR
operating ranges.
\item \textbf{Planning validation.} Compare weak-node, weak-link, conversion,
and coordinated-retrofitting rankings against full recomputation and against
SCR, impedance, Fiedler, flow loading, damping participation, and published
damping-index baselines.
\end{enumerate}

\section{What Is Actually New}

The safest novelty statement is:

\begin{quote}
The project formulates detailed IBR small-signal damping margin analysis as a
controller aware rational self-energy problem, validates the Schur NEP bridge
against full ANDES eigensolves including a control-family critical mode, and
uses that bridge to separate network-family and controller-family mechanisms.
\end{quote}

The strongest current pieces are:
\begin{enumerate}[leftmargin=*]
\item \textbf{Validated rational bridge.} The Phase-0D result is precise,
non-circular, and full-ANDES-based.
\item \textbf{Control-family criticality.} Mix60's PLL-family 13.7417 Hz mode
is the anchor that explains why static and fixed-point bridges fail.
\item \textbf{Resonant Braess mechanism.} Strict but minority cases show line
reinforcement moving a network-family mode toward a PLL-dominated control pole.
\item \textbf{Certificate discrepancy.} $S_\zeta$ and $S_g$ disagree under
several Gamma choices; $w+D$ and $w+M$ packages have protected-resolvent
coordination value. The mechanism remains open.
\item \textbf{Unified hierarchy.} The hierarchy itself is a synthesis, not a
single theorem. Its value is that it prevents overclaim: each broken classical
guarantee maps to a specific planning risk.
\end{enumerate}

\section{Submission Map}

\subsection*{Paper 1: Controller-aware rational bridge}
Core content: Level 2 bridge, Beyn contour validation, Mix60 control-family
mode, resonant Braess mechanism, scalar/fixed-point failures, and negative
controls. This is the nearest to a defensible submission.

\subsection*{Paper 2: Active nonviscous IBR grids}
Core content: damping-operator hierarchy and active-material framing. This
requires the passivity-defect gate and an adversarial literature search before
it can be claimed.

\subsection*{Paper 3: Flow-feasible protected-resolvent retrofitting}
Core content: bi-orthogonal resolvent sensitivity, spatial decoupling,
flow-feasible retrofit LP, and Gamma-robust certificate discrepancy. This needs
local mechanism descriptors and stronger comparisons against H2, SCR,
impedance, Fiedler, and transmission-switching baselines.

\section*{Acknowledgment of Dead Ends}

The project produced several attractive hypotheses that did not survive gates.
They should remain in the internal ledger because they prevent future
backsliding:
\begin{itemize}[leftmargin=*]
\item Nonlinear three-lever synergy as a headline: killed by Gate 1.
\item Universal inertia Braess: killed by fixed-ratio and Monte Carlo controls.
\item Second order always beating large reduced QEP: killed by cross-checks.
\item Inertia-control cross term as frequent ANDES mechanism: killed by 0/13
measurable cases.
\item Virtual-inertia placement over the controller-aware bridge: blocked by
missing physical GFM inertia parameter in REGF1 sheets.
\item Global non-normality or global commutator defect as explanation for
certificate discrepancy: weak in Gamma gate, not killed as a family but not
currently supported.
\end{itemize}

\begin{thebibliography}{99}

\bibitem{rayleigh1894}
J. W. Strutt (Lord Rayleigh), \emph{The Theory of Sound}, 2nd ed. London:
Macmillan, 1894.

\bibitem{caughey1965}
T. K. Caughey and M. E. J. O'Kelly, ``Classical normal modes in damped linear
dynamic systems,'' \emph{Journal of Applied Mechanics}, vol. 32, no. 3,
pp. 583--588, 1965.

\bibitem{biot1955}
M. A. Biot, ``Variational principles in irreversible thermodynamics with
application to viscoelasticity,'' \emph{Physical Review}, vol. 97, no. 6,
pp. 1463--1469, 1955.

\bibitem{golla1985}
D. F. Golla and P. C. Hughes, ``Dynamics of viscoelastic structures--a time
domain, finite element formulation,'' \emph{Journal of Applied Mechanics},
vol. 52, no. 4, pp. 897--906, 1985.

\bibitem{mctavish1993}
D. J. McTavish and P. C. Hughes, ``Modeling of linear viscoelastic space
structures,'' \emph{Journal of Vibration and Acoustics}, vol. 115, no. 1,
pp. 103--110, 1993.

\bibitem{adhikari2001}
S. Adhikari, ``Dynamics of non-viscously damped linear systems,''
\emph{Journal of Engineering Mechanics}, vol. 127, no. 3, pp. 328--339, 2001.

\bibitem{tisseur2001}
F. Tisseur and K. Meerbergen, ``The quadratic eigenvalue problem,''
\emph{SIAM Review}, vol. 43, no. 2, pp. 235--286, 2001.

\bibitem{beyn2012}
W.-J. Beyn, ``An integral method for solving nonlinear eigenvalue problems,''
\emph{Linear Algebra and its Applications}, vol. 436, no. 10, pp. 3839--3863,
2012.

\bibitem{trefethen2005}
L. N. Trefethen and M. Embree, \emph{Spectra and Pseudospectra: The Behavior of
Nonnormal Matrices and Operators}. Princeton, NJ: Princeton University Press,
2005.

\bibitem{kundur1994}
P. Kundur, \emph{Power System Stability and Control}. New York: McGraw-Hill,
1994.

\bibitem{rogers2000}
G. Rogers, \emph{Power System Oscillations}. Boston: Kluwer Academic
Publishers, 2000.

\bibitem{dorfler2014}
F. D\"orfler, M. Chertkov, and F. Bullo, ``Synchronization in complex networks
of phase oscillators: A survey,'' \emph{Automatica}, vol. 50, no. 6,
pp. 1539--1564, 2014.

\bibitem{poolla2017}
B. K. Poolla, S. Bolognani, and F. D\"orfler, ``Optimal placement of virtual
inertia in power grids,'' \emph{IEEE Transactions on Automatic Control},
vol. 62, no. 12, pp. 6209--6220, 2017.

\bibitem{pagnier2019}
L. Pagnier and P. Jacquod, ``Optimal placement of inertia and primary control
in high voltage power grids,'' \emph{arXiv:1906.06922}, 2019.

\bibitem{harnefors2007}
L. Harnefors, ``Modeling of three-phase dynamic systems using complex transfer
functions and transfer matrices,'' \emph{IEEE Transactions on Industrial
Electronics}, vol. 54, no. 4, pp. 2239--2248, 2007.

\bibitem{sun2011}
J. Sun, ``Impedance-based stability criterion for grid-connected inverters,''
\emph{IEEE Transactions on Power Electronics}, vol. 26, no. 11,
pp. 3075--3078, 2011.

\bibitem{witthaut2012}
D. Witthaut and M. Timme, ``Braess's paradox in oscillator networks,
desynchronization and power outage,'' \emph{New Journal of Physics}, vol. 14,
083036, 2012.

\bibitem{han2020}
Y. Han and D. J. Hill, ``Transmission switching for power system stability:
H2-performance and graph-theoretic viewpoints,'' [VERIFY exact bibliographic
details], 2020.

\bibitem{cui2021}
H. Cui, F. Li, and K. Tomsovic, ``Hybrid symbolic-numeric framework for power
system modeling and analysis,'' \emph{IEEE Transactions on Power Systems},
vol. 36, no. 2, pp. 1373--1384, 2021.

\bibitem{pan2020}
D. Pan, X. Fu, Q. Chen, P. Lu, and J. Tan, ``A modal perturbation method for
eigenvalue problem of non-proportionally damped system,'' \emph{Applied
Sciences}, vol. 10, no. 1, 341, 2020.

\bibitem{liyanage2025}
C. Liyanage \emph{et al.}, ``Strategic placement of grid-forming inverters
considering spatiotemporal dynamics and composite stability index,''
\emph{IEEE Open Journal of the Industrial Electronics Society}, vol. 6,
pp. 290--308, 2025.

\bibitem{setareh2025}
M. Setareh and M. S. Ghazizadeh, ``A systematic approach to compute power
system oscillatory modes damping-ratio sensitivities considering different
synchronous generator models,'' \emph{Computers and Electrical Engineering},
vol. 126, 110501, 2025.

\end{thebibliography}

\end{document}
